\documentclass[a4j,dvipdfmx]{jsarticle}
\usepackage{amsmath,amssymb}
\usepackage{siunitx}
\usepackage{ascmac}
\usepackage{amsthm}
\usepackage{bm}

\usepackage[margin=10truemm]{geometry}

\newcommand{\R}{\mathbb{R}}
\newcommand{\rank}{\mathrm{rank}}
\newcommand{\tr}{\mathrm{tr}}
\newcommand{\Ker}{\mathrm{Ker}\hspace{0.5mm}}
\newcommand{\E}{\mathcal{E}}
\newcommand{\F}{\mathcal{F}}
\renewcommand{\labelenumi}{(\arabic{enumi})}
\renewcommand{\Im}{\mathrm{Im}\hspace{0.5mm}}
% \usepackage[dvipdfmx]{hyperref}
\begin{document}
    \part*{テスト}
    \begin{align*}
        &y\in\mathrm{Im} f ならあるx\in Vがあってy=f(x)とかけるが, このxはVの基底e_1,\dots,e_nを用いてx=\sum x_{i}e_{i}とかける. \\ &よって, y=\sum x_{i}f(e_{i})とかけるから, yはf(e_1),\dots,f(e_n)の一次結合. この中から一次独立なものを選べばそれは基底になる.
    \end{align*}
\end{document}